\documentclass[12pt,a4paper]{article}

\usepackage[utf8]{inputenc}
\usepackage[T1]{fontenc}
\usepackage[spanish]{babel}
\usepackage{lmodern}
\usepackage{microtype}
\usepackage{geometry}
\usepackage{setspace}
\usepackage{parskip}

\geometry{
  top=2.5cm,
  bottom=2.5cm,
  left=3cm,
  right=3cm
}

\setstretch{1.15}

\title{Del interés, la importancia y la piedad}
\author{Vicente Domínguez Arca}
\date{}

\begin{document}

\maketitle

No es, en absoluto, pretensión de este autor convertir esta nota en una discusión sobre el fundamento del Derecho.

Pero me voy a permitir, como no experto, alguna licencia.

Tal vez el Derecho romano pueda considerarse no solamente una de las raíces, sino parte del tronco sobre el que posteriormente se desarrollaron muchos de los complejos Estados de derecho que encontramos, sobre todo, en sociedades occidentales u occidentalizadas.

No necesito que el lector asuma esta afirmación en toda su extensión.

Me basta con que conceda algo bastante más modesto: que el Derecho romano ha desempeñado un papel central en la construcción de numerosas categorías jurídicas posteriores.

Y entre esas categorías aparece una que resulta especialmente interesante para lo que aquí quiero plantear.

El \emph{interés}.

Conviene precisar inmediatamente que no pretendo reducir toda la complejidad de los usos jurídicos romanos del término a una única definición.

Me interesa particularmente la expresión \emph{id quod interest}.

Aquello que, en determinados contextos jurídicos, permitía atender a la diferencia entre la situación patrimonial efectiva de una persona y aquella en la que se habría encontrado si una obligación se hubiera cumplido o si determinado perjuicio no se hubiera producido.

Dicho de manera deliberadamente sencilla: la diferencia entre lo que se tiene y aquello que razonablemente se habría tenido de no haber ocurrido un determinado \emph{algo}.

Un incumplimiento.

Una pérdida.

Un daño.

Una privación.

Una interferencia.

Algo relativamente próximo, aunque naturalmente no idéntico en todos sus desarrollos, a aquello que hoy podemos pensar mediante determinadas formas de indemnización.

Y considero crucial detenernos aquí.

Porque el interés financiero contemporáneo parece conservar la palabra mientras ha desplazado considerablemente aquello que debe representar ese \emph{algo}.

Pensemos en un préstamo.

Una entidad financiera concede una determinada cantidad a un deudor y pasa, por ello, a convertirse en acreedora.

A partir de ese momento, sobre el principal aparece un interés.

La pregunta habitual consiste en determinar cuánto.

Qué porcentaje.

Qué plazo.

Qué riesgo.

Qué referencia.

Qué tipo de mercado.

Qué expectativas.

Qué objetivos comerciales.

Qué inflación esperada.

Qué perfil presenta el deudor.

Todo ello puede combinarse mediante mecanismos enormemente complejos hasta producir una cifra.

Pero la pregunta que interesa aquí no es exactamente cómo se calcula.

Es anterior.

¿Qué representa?

Si el interés fuera únicamente aquello que el acreedor deja de obtener por haber cedido temporalmente una riqueza que previamente poseía, la cuestión podría parecer relativamente intuitiva.

Quien dispone de algo y renuncia durante un tiempo a utilizarlo reclama una compensación por aquella privación.

Pero el crédito financiero contemporáneo complica considerablemente esa imagen.

Una entidad financiera no necesita disponer previamente, en forma de dinero acumulado y completamente separado de otros usos, de toda la cantidad que concede mediante un préstamo.

La propia concesión de crédito puede generar dinero bancario mediante la creación del depósito correspondiente, dentro, naturalmente, de las restricciones de capital, liquidez, regulación, riesgo y funcionamiento del sistema financiero.

No se trata de que una entidad pueda crear riqueza sin límite.

Tampoco de negar los costes y riesgos que asume.

Se trata simplemente de reconocer una particularidad fundamental:

el principal prestado no tiene por qué corresponderse con una cantidad idéntica de dinero previamente poseída por el acreedor y de cuyo uso este haya tenido que privarse.

Y entonces la pregunta cambia de intensidad.

Si aquello prestado puede surgir, al menos parcialmente, en el propio acto crediticio, ¿qué representa exactamente el interés exigido a cambio?

Más aún:

si el principal puede aparecer mediante la creación de dinero bancario, mientras que para el deudor queda constituida una obligación efectiva de devolver principal e intereses, ¿qué es exactamente aquello que estamos valorando cuando asignamos un precio al interés?

Obviamente existen respuestas económicas.

Riesgo.

Tiempo.

Liquidez.

Costes operativos.

Coste de financiación.

Expectativas de inflación.

Probabilidad de impago.

Beneficio.

Competencia.

Política monetaria.

Y una larga lista de variables que este autor no pretende aquí desarrollar.

Mi gusto podría llevarme ahora por ese derrotero y convertir esta nota en una extensa reflexión económica o macroeconómica sobre el funcionamiento financiero contemporáneo.

No lo haré.

Porque perderíamos el fundamento filosófico de la pregunta.

Lo que me interesa es algo bastante más sencillo:

sabemos calcular el interés.

Pero ¿sabemos exactamente qué estamos llamando interés?

Y es aquí donde la propia palabra empieza a resultar sugerente.

Más allá de sus usos jurídicos particulares, \emph{interesse} contiene una estructura lingüística enormemente expresiva.

\emph{Inter}.

\emph{Esse}.

Estar entre.

Ser entre.

No pretendo reducir toda la evolución semántica de la palabra a su composición latina.

Las palabras evolucionan.

Cambian.

Adquieren sentidos que no se encuentran contenidos de manera mecánica en sus raíces.

Pero las raíces pueden servir, una vez más, como brújula.

Y resulta cuando menos curioso que aquello que hoy llamamos \emph{interés} se encuentre relacionado, en nuestra experiencia lingüística cotidiana, con aquello que nos concierne, aquello que reclama nuestra atención, aquello que se sitúa de algún modo entre nosotros y otra cosa.

Algo nos interesa.

Algo tiene interés.

Algo despierta nuestro interés.

Y entonces aparece otra palabra cercana, aunque no proceda exactamente de la misma genealogía:

\emph{importante}.

\emph{Importare}.

Traer dentro.

Llevar hacia dentro.

Aquello que merece ser traído hacia nuestra consideración.

Aquello que posee suficiente peso como para ser incorporado a nuestro espacio de atención.

La relación entre ambas palabras no es etimológicamente directa.

Pero semánticamente resulta difícil ignorarla.

Aquello que nos interesa suele ser aquello que nos importa.

Y aquello que nos importa suele adquirir interés.

No se trata de un simple comentario filológico.

Al menos no para este autor.

Porque quizá la operativa contemporánea del interés financiero, aunque no se corresponda sin más con aquel \emph{id quod interest} romano, sí termine acercándose de un modo sorprendente a aquello que resulta interesante o importante para el acreedor.

El interés deja entonces de aparecer exclusivamente como una medida de aquello que se ha perdido.

Puede aparecer también como una medida de aquello que se espera obtener.

Y aquí surge una expresión deliberadamente ambigua:

el interés sobre el interés.

El interés económico determinado por aquello que interesa al acreedor.

Aquello que considera importante.

Aquello que necesita obtener.

Aquello que pretende ganar.

Aquello que el mercado permite exigir.

Aquello que el riesgo parece justificar.

Aquello que las expectativas convierten en razonable.

No afirmo que estas dimensiones conviertan automáticamente al interés en algo ilegítimo.

Tampoco pretendo reducir la relación acreedor-deudor a una escena elemental de dominación.

Sería demasiado sencillo.

Lo que intento señalar es otra cosa.

La palabra parece conservar una apariencia de continuidad mientras aquello que designa se ha vuelto mucho más complejo.

Y quizá esa continuidad lingüística nos impida, en ocasiones, volver a preguntar por la naturaleza de aquello que damos por supuesto.

Decimos interés.

Calculamos interés.

Pagamos interés.

Cobramos interés.

Negociamos interés.

Pero quizá pronunciamos demasiado rápidamente una palabra cuyo contenido damos por completamente resuelto.

Y aquí aparece, finalmente, la piedad.

No utilizo la palabra únicamente en su sentido contemporáneo de compasión.

Me interesa también su raíz latina, \emph{pietas}, vinculada al deber, a la consideración debida, a la relación responsable con aquello a lo que uno se encuentra ligado.

Quizá deberíamos tener cierta piedad con nuestras palabras.

No porque debamos obedecerlas.

No porque su primer significado tenga derecho a gobernar todos los significados posteriores.

Las lenguas no funcionan así.

Pero sí porque nos debemos, en cierta medida, a ellas.

No me refiero a las palabras pronunciadas por uno mismo.

Me refiero a las palabras que heredamos y mediante las cuales podemos comunicarnos.

Pensamos con ellas.

Discutimos con ellas.

Construimos instituciones con ellas.

Escribimos leyes con ellas.

Establecemos contratos con ellas.

Y, en ocasiones, naturalizamos realidades enteras mediante ellas.

Por eso quizá la precisión lingüística no sea únicamente una cuestión estética.

Puede ser también una cuestión de responsabilidad.

Cuando una institución conserva una palabra mientras transforma profundamente aquello que esa palabra designa, quizá convenga detenerse y preguntar qué continuidad estamos aceptando sin examinarla.

No necesariamente para regresar a un significado antiguo.

Mucho menos para convertir la etimología en una autoridad.

Sino para observar el desplazamiento.

Para saber qué hemos añadido.

Qué hemos eliminado.

Qué hemos transformado.

Y qué hemos terminado dando por evidente.

Quizá, en la perspectiva de este autor, algunos de los problemas que hoy parecen extraordinariamente complejos podrían comenzar a pensarse con mayor claridad si prestáramos algo más de atención a las palabras mediante las que los formulamos.

No porque las palabras resuelvan los problemas.

Pero sí porque pueden ocultarlos.

Tal vez debamos preguntarnos con mayor frecuencia no solamente cuánto interés corresponde.

Sino qué es exactamente aquello que estamos llamando interés.

Qué está entre.

Qué nos importa.

Y, sobre todo, a quién.

Quizá sea solamente una cuestión de claridad.

O quizá, en cierto sentido, también de piedad.

\end{document}
